% =====================================================================
%  CAPÍTULO 4 — METODOLOGIA E PROCESSO DE DESENVOLVIMENTO
% =====================================================================

\chapter[METODOLOGIA E PROCESSO DE DESENVOLVIMENTO]{METODOLOGIA E PROCESSO DE\\DESENVOLVIMENTO}
\label{cap:metodologia}

Este capítulo descreve como o trabalho foi conduzido. Ele não se limita a
listar ferramentas: explicita o método de pesquisa, o processo de engenharia, a
linha de base registrada no INPI, os dados coletados e o desenho da avaliação.
A intenção é permitir que outro pesquisador entenda o que foi feito, por que foi
feito e o que pode ser reproduzido.

\section{Design Science Research}
\label{sec:dsr}

O trabalho adota a \textit{Design Science Research} (DSR) como abordagem
metodológica. A DSR é indicada quando o objeto de pesquisa é um artefato
projetado para resolver um problema prático e quando o conhecimento é produzido
pela construção e pela avaliação desse artefato \cite{hevner2004}. O artefato,
neste caso, é o \argus{}: um sistema de software, mas também um conjunto de
decisões de arquitetura, modelos, protocolos e documentação.

Hevner et al. \cite{hevner2004} organizam a DSR em torno de sete diretrizes:
projeto como artefato, relevância do problema, avaliação, contribuição de
pesquisa, rigor metodológico, processo de projeto e comunicação. Peffers et al.
\cite{peffers2007} propõem um processo de seis etapas que torna o ciclo mais
explícito. A Tabela~\ref{tab:dsr} relaciona essas etapas ao presente trabalho.

\begin{table}[htbp]
    \centering
    \caption{Etapas da \textit{Design Science Research} aplicadas ao trabalho}
    \label{tab:dsr}
    \small
    \begin{tabular}{
        >{\raggedright\arraybackslash}p{3.4cm}
        >{\raggedright\arraybackslash}p{7.2cm}
        >{\raggedright\arraybackslash}p{4.0cm}
    }
        \toprule
        \textbf{Etapa} & \textbf{O que foi feito} & \textbf{Onde aparece} \\
        \midrule
        Identificação do problema &
        Caracterização da fragmentação entre percurso, simulação e análise;
        definição dos cinco problemas e dos requisitos &
        Capítulos~\ref{cap:introducao} e~\ref{cap:estado-arte} \\
        Objetivos da solução &
        Definição do objetivo geral e dos objetivos específicos, com
        delimitação de escopo &
        Capítulo~\ref{cap:introducao} \\
        Projeto e desenvolvimento &
        Arquitetura, escolha de tecnologias, implementação dos módulos e
        ciclos iterativos &
        Capítulos~\ref{cap:arquitetura} e~\ref{cap:implementacao} \\
        Demonstração &
        Execução do fluxo completo no desktop e de sessões ao vivo no Android,
        nos cenários de reconhecimento, exploração e treino &
        Capítulo~\ref{cap:resultados} \\
        Avaliação &
        Estudo piloto com o autor, testes técnicos, verificação de arquivos
        FIT e registro de problemas de uso &
        Capítulo~\ref{cap:resultados} \\
        Comunicação &
        Redação do TCC, documentação do repositório, registro no INPI e
        disponibilização do código &
        Capítulo~\ref{cap:conclusao} e Apêndice~\ref{ap:registro-inpi} \\
        \bottomrule
    \end{tabular}
    \vspace{0.2cm}
    \par\small Fonte: adaptado de Peffers et al. \cite{peffers2007}.
\end{table}

A DSR não elimina a necessidade de rigor experimental; ela desloca o foco para a
adequação entre problema, artefato e avaliação. Por isso, o capítulo distingue
três tipos de evidência: a \textbf{evidência de construção} (o código e o
histórico de desenvolvimento), a \textbf{evidência de demonstração} (o fluxo
executado em cenários reais) e a \textbf{evidência de avaliação} (o piloto e os
testes técnicos). Nenhuma delas é tratada como prova de eficácia clínica ou de
generalização estatística.

\section{Problema, requisitos e escopo}
\label{sec:problema-requisitos}

A Tabela~\ref{tab:requisitos} sintetiza os requisitos que orientaram o
desenvolvimento. Eles foram derivados dos usos descritos no
Capítulo~\ref{cap:introducao} e refinados ao longo dos ciclos de implementação.

\begin{table}[htbp]
    \centering
    \caption{Requisitos funcionais e não funcionais do Argus}
    \label{tab:requisitos}
    \small
    \begin{tabular}{
        >{\raggedright\arraybackslash}p{1.4cm}
        >{\raggedright\arraybackslash}p{7.4cm}
        >{\raggedright\arraybackslash}p{5.4cm}
    }
        \toprule
        \textbf{ID} & \textbf{Requisito} & \textbf{Critério de verificação} \\
        \midrule
        RF01 & Importar GPX e extrair distância, altitude e inclinação &
        Perfil altimétrico e distância conferem com o arquivo de origem \\
        RF02 & Controlar o rolo pela inclinação do percurso &
        No desktop, resistência varia ao longo da rota; modo simulação estável \\
        RF03 & Registrar e exibir telemetria de potência, FC e cadência &
        Desktop grava amostras por notificação BLE; HUD atualizado nas duas
        plataformas \\
        RF04 & Calcular NP, IF, TSS, TRIMP, PMC, HRR e desacoplamento &
        No desktop, fórmulas documentadas e valores conferidos em casos de
        teste \\
        RF05 & Persistir perfis, atividades e recordes localmente &
        No desktop, dados permanecem após reinício e perfis ficam isolados \\
        RF06 & Exportar FIT e permitir integrações opcionais &
        No desktop, arquivo aceito por Strava/Intervals.icu; integração pode ser
        desligada \\
        RF07 & Executar em desktop e Android &
        Fluxo completo no desktop; sessão ao vivo no Android, com adaptações de
        UI \\
        RF08 & Apoiar o fluxo manual de reconhecimento &
        No desktop, GPX, perfil, mapa, sessão e análise do FIT; no Android,
        GPX, perfil, mapa e sessão ao vivo \\
        RNF01 & Processar localmente, sem nuvem obrigatória &
        No desktop, cálculo e histórico funcionam sem conta; Android mantém a
        sessão local; rede só para mapas/integrações \\
        RNF02 & Interface responsiva e utilizável em telas pequenas &
        Layout adaptado; navegação testada em telefone Android \\
        RNF03 & Código aberto e licenciado sob GPLv3 &
        Repositório com licença; código disponível para auditoria \\
        RNF04 & Preservar integridade e recuperação de sessões &
        No desktop, exportação FIT válida e recuperação de sessão interrompida \\
        RNF05 & Explicitar fórmulas, parâmetros e limitações &
        Documentação e interface indicam o que é medido ou estimado \\
        \bottomrule
    \end{tabular}
    \vspace{0.2cm}
    \par\small Fonte: autoria própria.
\end{table}

O escopo foi deliberadamente limitado. O \argus{} não inclui catálogo de rotas,
reconhecimento automático de curvas, análise de vídeo, competição massiva,
\textit{multiplayer} em tempo real ou recomendações clínicas. Essas funções
poderiam ser úteis, mas ampliariam o escopo para além do que um TCC consegue
implementar e avaliar com rigor. A Seção~\ref{sec:delimitacoes} do
Capítulo~\ref{cap:introducao} explicita esses limites.

\section{Processo de desenvolvimento}
\label{sec:processo-desenvolvimento}

O desenvolvimento seguiu um processo iterativo e incremental, com ciclos curtos
de implementação, uso e correção. O histórico Git do repositório registra 301
\textit{commits}, 130 \textit{merges} e 171 referências a \textit{issues}, entre
6 de janeiro e 18 de outubro de 2026 (levantamento realizado em 18/10/2026).
Esse histórico não é apresentado como métrica de qualidade, mas como evidência
de continuidade e de rastreabilidade das decisões.

O processo pode ser descrito em três ciclos principais.

\subsection{Primeiro ciclo: núcleo de simulação e linha de base}
\label{sec:primeiro-ciclo}

O primeiro ciclo construiu o núcleo funcional: conexão BLE com sensor de
frequência cardíaca e rolo, leitura de GPX, cálculo de inclinação, mapa 3D,
persistência local e exportação FIT. Esse ciclo produziu a versão registrada no
INPI em 20/01/2026
(Apêndice~\ref{ap:registro-inpi}). A linha de base já continha o que o trabalho
considera essencial: o percurso entrava, o rolo respondia e a sessão era
registrada.

\subsection{Segundo ciclo: análise, treino e integrações}
\label{sec:segundo-ciclo}

O segundo ciclo ampliou o sistema para além do simulador. Foram implementados
treinos estruturados em ZWO, testes de FTP e \textit{ramp test}, modos de
evento (sprint, KOM e contrarrelógio), análise de PMC, TRIMP, HRR e
desacoplamento cardiovascular, além das integrações com Strava e Intervals.icu.
O foco passou a ser a interpretação da sessão, e não apenas a sua execução.

\subsection{Terceiro ciclo: mobilidade, mapa e experiência}
\label{sec:terceiro-ciclo}

O terceiro ciclo levou o sistema para Android por meio do Capacitor, migrou o
mapa para MapLibre GL JS com terreno, unificou os modos de estúdio e desafio,
adicionou o assistente de IA local via Ollama, o rastreamento de desgaste de
componentes e melhorias de estabilidade na camada BLE. Esse ciclo também
produziu a maior parte das decisões de interface responsiva que hoje sustentam
o modo móvel. No Android, a entrega foi deliberadamente menor: sessão ao vivo,
GPX, ZWO, BLE/FE-C e HUD; FIT, histórico, testes de aptidão e integrações
permaneceram no desktop.

\subsection{Evidência de evolução}
\label{sec:evolucao}

A Tabela~\ref{tab:evolucao} compara a versão registrada com a versão
apresentada neste trabalho.

\begin{table}[htbp]
    \centering
    \caption{Evolução do Argus: versão registrada \textit{versus} versão atual}
    \label{tab:evolucao}
    \small
    \begin{tabular}{
        >{\raggedright\arraybackslash}p{3.2cm}
        >{\raggedright\arraybackslash}p{5.6cm}
        >{\raggedright\arraybackslash}p{5.6cm}
    }
        \toprule
        \textbf{Aspecto} & \textbf{Versão registrada (01/2026)} &
        \textbf{Versão atual (2026)} \\
        \midrule
        Plataforma & Desktop (Wails) & Desktop + modo Android (Capacitor) \\
        Motor de mapa & Mapbox 3D & MapLibre GL JS (Terrain-RGB) \\
        Motor físico & Go (\textit{backend}) & Go no desktop; reimplementado em JavaScript no Android \\
        Conectividade BLE & FTMS + HRM & Desktop: FTMS, FE-C, HRM, potência e rolo virtual; Android: FE-C, potência e HRM \\
        Treinos & Livres & Livres + estruturados (ZWO) + testes FTP \\
        Métricas & Básicas & NP, IF, TSS, TRIMP, PMC, HRR, desacoplamento \\
        Integrações & Não havia & Strava (OAuth 2.0) e Intervals.icu \\
        IA & Não havia & Assistente local via Ollama (planos ZWO) \\
        Gamificação & Não havia & XP, níveis, troféus, conquistas, metas \\
        Modos de treino & Não havia & Estúdio, Sprint, KOM, Contrarrelógio \\
        Manutenção & Não havia & Rastreamento de desgaste de componentes \\
        Persistência & Perfil único & Multiperfil isolado no desktop (SQLite); Android sem histórico \\
        \bottomrule
    \end{tabular}
    \vspace{0.2cm}
    \par\small Fonte: autoria própria.
\end{table}

A evolução não representa abandono da versão registrada, mas maturação. No
desktop, as funcionalidades essenciais da linha de base (simulação BLE, mapa,
exportação FIT e histórico local) foram mantidas e ampliadas. O Android
acrescenta uma sessão ao vivo com GPX, ZWO e BLE, mas não exporta FIT nem mantém
histórico. O registro no INPI protege a base do software; este documento
formaliza o estado atual e o processo que levou de uma versão à outra.

\section{Demonstração}
\label{sec:demonstracao}

A demonstração foi organizada em três cenários, correspondentes aos usos
descritos no Capítulo~\ref{cap:introducao}. No desktop, os três foram
executados de ponta a ponta; no Android, o reconhecimento e a exploração foram
exercitados como sessão ao vivo, sem a etapa de FIT e análise posterior.

\begin{enumerate}
    \item \textbf{Reconhecimento}: importar o GPX de uma prova, inspecionar o
        perfil, pedalar o percurso e comparar a resposta fisiológica com a
        leitura prévia do terreno. O caso principal foi o percurso
        Geraardsbergen--Leberg (aproximadamente 21~km, 297~m de ganho e rampas
        de até 18,2\%), escolhido por combinar subidas curtas e íngremes com
        trechos de recuperação.
    \item \textbf{Exploração}: pedalar uma rota de outra região sem deslocamento,
        verificando se o mapa e o perfil permitem entender a geometria do lugar
        e se a experiência faz sentido para uso recreativo.
    \item \textbf{Treino}: executar uma sessão livre ou estruturada, exportar o
        FIT e analisar as métricas no próprio \argus{} e em uma plataforma
        externa.
\end{enumerate}

O repositório de trabalho contém 23 arquivos GPX, 5 arquivos ZWO e 274 arquivos
FIT, que documentam a variedade de percursos e sessões usados ao longo do
desenvolvimento. A demonstração não exige que todos sejam analisados; eles
formam o acervo de onde os casos foram selecionados. Os arquivos FIT são
evidência do fluxo de desktop; o Android não os produz nesta versão.

\section{Avaliação}
\label{sec:avaliacao}

\subsection{Desenho do estudo piloto}
\label{sec:desenho-piloto}

A avaliação principal é um \textbf{estudo piloto exploratório de caso único},
com o próprio autor como participante. O autor utiliza o sistema em sessões de
30 a 60 minutos, com rolo e sensor de frequência cardíaca reais, em
\textit{desktop} e em telefone Android. O objetivo é verificar a viabilidade do
fluxo, identificar
falhas de integração e levantar hipóteses para estudos futuros, não estimar
efeitos de treinamento.

O piloto combina dados quantitativos e qualitativos:

\begin{itemize}
    \item \textbf{conclusão do fluxo}: no desktop, importação do GPX, conexão
        dos sensores, início da sessão, controle do rolo, encerramento e
        exportação; no Android, importação, conexão e sessão ao vivo;
    \item \textbf{estabilidade}: no desktop, registro de desconexões, falhas de
        escrita, perda de amostras e recuperação de sessão;
    \item \textbf{integridade do FIT}: no desktop, aceitação do arquivo por
        Strava e Intervals.icu e conferência de distância, duração, potência e
        elevação;
    \item \textbf{usabilidade}: diário de uso com problemas observados,
        dificuldades de navegação e ajustes necessários;
    \item \textbf{desempenho}: comportamento da interface e do mapa durante a
        sessão, sem transformar a percepção visual em medida de latência.
\end{itemize}

\subsection{Análise dos dados}
\label{sec:analise-dados}

Os dados quantitativos são apresentados de forma descritiva: contagens,
durações, valores de potência e frequência cardíaca, taxas de sucesso e
incidência de falhas. Não há teste de hipótese nem inferência estatística,
porque a amostra é única e não representa uma população. Os dados qualitativos
são organizados por tema (conexão, mapa, controle, análise, mobilidade) e
usados para explicar os resultados quantitativos.

\subsection{Validade e limitações}
\label{sec:validade}

No \textit{desktop}, a validade interna do piloto é reforçada pela triangulação
entre o registro da sessão, o arquivo FIT exportado e a observação do autor. A validade externa é
baixa: um único participante, um conjunto restrito de equipamentos e um
contexto específico de uso. Os resultados não permitem afirmar que o sistema é
usável por outros ciclistas, que melhora o desempenho ou que substitui um
treinador. Permitem afirmar que o fluxo funciona nas condições testadas e que
os problemas identificados são reais.

Também não foram realizados testes com outros participantes, questionários
padronizados de usabilidade, avaliação clínica de VO\textsubscript{2}max ou
medição de lactato. Essas ausências são limitações declaradas e devem orientar
trabalhos futuros.

\section{Aspectos éticos}
\label{sec:etica}

O estudo utiliza apenas dados do próprio autor. Não há participantes humanos
externos, coleta de dados de terceiros, material biológico ou intervenção
clínica. Os dados de treino são armazenados localmente e podem ser apagados pelo
próprio usuário. As integrações com Strava e Intervals.icu são opcionais e
dependem de autorização explícita do usuário.

A pesquisa segue os princípios da Lei Geral de Proteção de Dados (LGPD) no que
se aplica ao tratamento de dados pessoais: minimização, finalidade e controle
do titular. Como o único titular é o autor, o risco é reduzido, mas o desenho do
sistema foi mantido compatível com o uso por terceiros em trabalhos futuros.

\section{Síntese do capítulo}
\label{sec:sintese-metodologia}

A metodologia combina DSR, desenvolvimento iterativo e avaliação piloto. A DSR
fornece a estrutura de problema, artefato e avaliação; o desenvolvimento
iterativo fornece a evidência de construção; o piloto fornece a evidência de
uso. O próximo capítulo detalha a arquitetura que sustenta o artefato e as decisões
que permitiram executar o fluxo completo no \textit{desktop} e uma sessão ao
vivo no Android, com o \textit{frontend} compartilhado.
